\documentclass[12pt,letterpaper]{article}
\usepackage[utf8]{inputenc}
\usepackage[spanish]{babel}
\usepackage[margin=2cm]{geometry}
\usepackage{enumitem}
\usepackage{hyperref}

\pagestyle{empty}

\begin{document}

% ==================== PAGINA 1: PREGUNTAS ====================
\begin{center}
    {\Large\bfseries Cuestionario -- Computaci\'on en la Nube}\\[0.3cm]
    {\normalsize Sistemas Distribuidos -- Grupo 01 -- Semestre 2026-2}\\[0.3cm]
    {\normalsize Equipo: García Cortés Adolfo de Jesus.}\\[0.3cm]
    {\normalsize Lara Hernandez Angel Husiel.}\\[0.3cm]
    {\normalsize Lugo Manzano Rodrigo.}\\[0.3cm]
    {\normalsize\url{https://drive.google.com/file/d/1fRmsMqElKWidYCqTxGynC-XjsWI1CXae/view?usp=sharing}}
\end{center}

\vspace{0.5cm}

\noindent\textbf{Instrucciones:} Seleccione la opci\'on correcta para cada pregunta.

\vspace{0.5cm}

\begin{enumerate}[label=\textbf{\arabic*.}, leftmargin=1.5em, itemsep=0.8cm]

    \item \textbf{¿Qué algoritmos de consenso distribuido entran en acción para mantener sincronizados los equipos y garantizar la consistencia absoluta de los datos ante cualquier fallo?}
    \begin{enumerate}[label=\alph*), nosep]
        \item Algoritmos de enrutamiento como Dijkstra y Bellman-Ford
        \item Algoritmos de consenso como Paxos y Raft
        \item Algoritmos de cifrado asimétrico como RSA y ECC
        \item Algoritmos de procesamiento masivo como MapReduce
    \end{enumerate}

    \item \textbf{?`Qu\'e tecnolog\'ia, es el est\'andar de factor para la orquestaci\'on de contenedores en la nube?}

    \begin{enumerate}[label=\alph*), nosep]
        \item Docker
        \item OpenStack
        \item Kubernetes
        \item VMware vSphere
    \end{enumerate}

    \item \textbf{?`C\'omo se denomina el mecanismo que agrega m\'as nodos al cl\'uster en lugar de aumentar la capacidad de un solo servidor para resolver la escalabilidad?}

    \begin{enumerate}[label=\alph*), nosep]
        \item Escalamiento vertical
        \item Escalamiento horizontal
        \item Balanceo de carga est\'atico
        \item Migraci\'on en caliente
    \end{enumerate}

    \item \textbf{?`Qu\'e empresa migr\'o completamente su infraestructura a Amazon Web Services tras una corrupci\'on grave en su base de datos en 2008, siendo el caso m\'as emblem\'atico de migraci\'on a la nube?}

    \begin{enumerate}[label=\alph*), nosep]
        \item Spotify
        \item Shopify
        \item Capital One
        \item Netflix
    \end{enumerate}

\end{enumerate}

\newpage

% ==================== PAGINA 2: RESPUESTAS ====================
\begin{center}
    {\Large\bfseries Respuestas del Cuestionario}\\[0.3cm]
    {\normalsize Computaci\'on en la Nube (Cloud Computing)}
\end{center}

\vspace{0.5cm}

\begin{enumerate}[label=\textbf{\arabic*.}, leftmargin=1.5em, itemsep=0.5cm]

    \item \textbf{Respuesta: b) Algoritmos de consenso como Paxos y Raft}

    La nube es un sistema distribuido a gran escala. Para que miles de computadoras independientes funcionen como un sistema único y mantengan los datos sincronizados, se utilizan algoritmos de consenso distribuido (como Paxos y Raft) que garantizan la consistencia aunque ocurran fallos en los nodos.
    \item \textbf{Respuesta: c) Kubernetes}

    Kubernetes fue desarrollado por Google y es mantenido por la Cloud Native Computing Foundation. Automatiza el despliegue, escalamiento y gesti\'on de contenedores.

    \item \textbf{Respuesta: b) Escalamiento horizontal}

    Consiste en agregar m\'as nodos al cl\'uster en lugar de aumentar la capacidad de un solo servidor. La nube aprovisiona servidores en segundos y los libera cuando la demanda baja.

    \item \textbf{Respuesta: d) Netflix}

    Tras una corrupci\'on en su base de datos en 2008, Netflix migr\'o completamente a AWS. El proceso tom\'o siete a\~nos y concluy\'o en 2016, adoptando microservicios y bases NoSQL.

\end{enumerate}

\end{document}
